\documentclass[aspectratio=169,UTF8]{ctexbeamer}
\usetheme{Madrid}
\usecolortheme{default}
\definecolor{DeepBlue}{HTML}{2563EB}
\definecolor{DeepInk}{HTML}{172033}
\definecolor{DeepGreen}{HTML}{059669}
\setbeamercolor{structure}{fg=DeepBlue}
\setbeamercolor{title}{fg=white,bg=DeepBlue}
\setbeamercolor{frametitle}{fg=white,bg=DeepBlue}
\setbeamercolor{alerted text}{fg=DeepGreen}
\setbeamertemplate{navigation symbols}{}
\setbeamertemplate{footline}[frame number]
\usepackage{booktabs}
\usepackage{listings}
\usepackage{xcolor}
\usepackage{hyperref}
\hypersetup{colorlinks=true,urlcolor=DeepBlue}
\lstset{
  basicstyle=\ttfamily\scriptsize,
  keywordstyle=\color{DeepBlue},
  stringstyle=\color{DeepGreen},
  commentstyle=\color{gray},
  showstringspaces=false,
  breaklines=true,
  frame=single,
  columns=fullflexible
}
\newcommand{\api}[1]{\texttt{#1}}
\newcommand{\keypoint}[1]{\alert{#1}}
\AtBeginSection[]{
  \begin{frame}{本节路线}
    \tableofcontents[currentsection]
  \end{frame}
}

\title{第8课：Files API}
\subtitle{上传、分页、引用、过期与删除}
\author{DeepSeek API 全栈课程}
\date{文档快照：2026-11-05}
\begin{document}
\frame{\titlepage}
\begin{frame}{学习目标}
\begin{itemize}
  \item 管理图片文件的完整生命周期
  \item 在 Chat 中复用 file\_id
  \item 掌握配额、分页、过期和账号归属
  \item 区分 OpenAI 与 Anthropic 兼容 Files API
\end{itemize}
\end{frame}
\section{生命周期}
\begin{frame}{何时用 Files API}
\begin{itemize}
  \item 同一图片要在多个请求中复用
  \item 内联请求会超过 48 MiB
  \item 图片超过 32 MiB、但不超过 64 MiB
  \item 希望上传与推理解耦
\end{itemize}
当前仅支持图片：JPEG、PNG、GIF、WebP。
\end{frame}
\begin{frame}{四个 OpenAI 兼容端点}
\begin{description}
  \item[POST /files] multipart 上传
  \item[GET /files] 分页列出
  \item[GET /files/\{id\}] 查询元数据
  \item[DELETE /files/\{id\}] 删除
\end{description}
file\_id 形如 \api{file-api-...}。
\end{frame}
\begin{frame}{上传字段}
\begin{itemize}
  \item \api{file}：必填，最大 64 MiB
  \item \api{purpose=user\_data}：必填
  \item \api{expires\_after[anchor]=created\_at}
  \item \api{expires\_after[seconds]}：3600--2592000
\end{itemize}
不传过期字段则永久有效；上传需在 10 分钟内完成。
\end{frame}
\begin{frame}[fragile]{上传与引用}
\begin{lstlisting}[language=Python]
with open("image.jpg", "rb") as handle:
    uploaded = client.files.create(
        file=handle, purpose="user_data")

content = [
  {"type": "text", "text": "分析图片"},
  {"type": "file", "file_id": uploaded.id},
]
\end{lstlisting}
文件属于账号，同账号所有 API Key 和 API 家族可共享。
\end{frame}
\section{列表与治理}
\begin{frame}{分页参数}
\begin{description}
  \item[after] file\_id 游标
  \item[limit] 1--1000
  \item[order] asc（默认）或 desc
  \item[purpose] 仅 user\_data
\end{description}
响应包含 data、first\_id、last\_id、has\_more。
\end{frame}
\begin{frame}{文件对象}
\begin{itemize}
  \item id、object、bytes、created\_at、filename、purpose
  \item 只有设置过期时才出现 expires\_at
  \item 文件名最长 512 字节
  \item 单用户最多 25 GiB、10000 个文件
\end{itemize}
上传成功后应持久化 file\_id 与业务资源的映射。
\end{frame}
\begin{frame}{清理策略}
\begin{itemize}
  \item 临时任务设置到期时间
  \item 业务删除时同步删除远端文件
  \item 定期分页审计孤儿文件与容量
  \item 删除是破坏性操作：先校验归属和引用计数
\end{itemize}
\end{frame}
\section{内容块与兼容}
\begin{frame}{file 块的两种来源}
\begin{itemize}
  \item \api{file\_id}：引用已上传文件
  \item \api{file\_data}：data URL 内联；可配 filename
\end{itemize}
file\_id 与 file\_data 互斥；filename 不能与 file\_id 同时出现。
\end{frame}
\begin{frame}{Anthropic 兼容差异}
\small
所有请求需 \api{anthropic-beta: files-api-2025-04-14}。裸 HTTP 路径为 \api{/anthropic/v1/files}。
\begin{itemize}
  \item 游标：after\_id / before\_id（互斥），默认 limit=20
  \item 无 order/purpose；大小字段为 size\_bytes
  \item type 替代 object；created\_at 为 RFC 3339
  \item 删除返回 type=file\_deleted
\end{itemize}
\end{frame}
\begin{frame}{失败与安全}
\begin{itemize}
  \item 上传前检查真实格式、大小与业务授权
  \item 不信任用户文件名；展示时转义
  \item 对 413 不重试原请求，应压缩或改方案
  \item file\_id 不是公开 URL，但仍属于敏感业务标识
\end{itemize}
\end{frame}
\begin{frame}{本课例程}
\begin{enumerate}
  \item \api{01\_file\_lifecycle.py}：上传/查询/列表/可选删除
  \item \api{02\_file\_vision.py}：file\_id 视觉问答
\end{enumerate}
删除需显式传 \api{--delete}，防止演示误删。
\end{frame}
\begin{frame}{官方资料}
\small
\url{https://api-docs.deepseek.com/zh-cn/guides/files_api}\\
\url{https://api-docs.deepseek.com/zh-cn/api/create-file}\\
\url{https://api-docs.deepseek.com/zh-cn/api/list-files}\\
\url{https://api-docs.deepseek.com/zh-cn/api/retrieve-file}\\
\url{https://api-docs.deepseek.com/zh-cn/api/delete-file}
\end{frame}
\end{document}
